% The command line: every medorama command and subcommand, the viewers' programs (vitalstream-rs,
% -cpp, -py; muse-py and dicom-py with their serve), and the repository's tools. The options are
% those of the code: medorama/src/main.rs (clap), viewers/vitalstream/rust/src/main.rs,
% cpp/vitalstream/src/main.cpp, python/vitalstream/vitalstream_py/__main__.py,
% viewers/ecg/python/muse/muse_reader_py/cli.py and page_server.py,
% viewers/pacs/python/dicom/dicom_index_py/cli.py, and each script of tools/.
\chapter{The command line}
\label{chap:cli}

This chapter is the reference of every program that Medorama and its viewers install: the
command \cmd{medorama}, which creates, fills and runs a catalog; the three VitalStream viewers,
\cmd{vitalstream-rs}, \cmd{vitalstream-cpp} and \cmd{vitalstream-py}; the database tools
\cmd{muse-py} and \cmd{dicom-py}, whose \cmd{serve} runs the ECG and the PACS viewer; and the
scripts of the repository's folder \file{tools/}. Each command is one entry: its syntax, what it
does, each option with its default, and an example. The tasks that the commands serve are
described in their own chapters: setting up and looking after a catalog in
\cref{chap:administration}, users and privacy in \cref{chap:users-privacy}, importing in
\cref{chap:importing}, recordings in \cref{chap:recordings}, add-ons in \cref{chap:addons}; the
settings files in \cref{chap:settings-reference}.

\section{Running the commands}
\label{sec:cli-running}

\cmd{medorama} is one program with subcommands: \cmd{medorama init}, \cmd{medorama serve},
\cmd{medorama import-table}, and so on. A subcommand that has its own subcommands is written in
two words (\cmd{medorama imports list}, \cmd{medorama user add}). Options are long options
(\opt{-{}-results DIR}); an option that takes a value is given it as the next argument or after
an equals sign (\opt{-{}-port 8090} or \opt{-{}-port=8090}); options and positional arguments may
come in any order after the subcommand.

\begin{params}
\item[-h, --help] Prints the usage of the program, or of a subcommand
  (\cmd{medorama import-table -{}-help}), and exits. \cmd{medorama help} \textit{COMMAND} does the
  same.
\item[-V, --version] Prints the version (\code{medorama 0.1.0}) and exits.
\item[--home DIR] The catalog's home folder, which holds \file{medorama.toml}, the SQLite
  database, the folders \file{actions/} and \file{formats/}. It is accepted by every subcommand,
  before or after it (\cmd{medorama -{}-home DIR serve}, \cmd{medorama serve -{}-home DIR}).
\end{params}

\paragraph{Which catalog a command uses.} A command uses the first of these that applies:
\begin{enumerate}
\item \opt{-{}-home DIR};
\item the environment variable \code{MEDORAMA\_HOME}, when it is set and not empty;
\item the catalog that \cmd{init} created last, or the one that \cmd{medorama -{}-home DIR use}
  chose, when that folder still holds a \file{medorama.toml}: its folder is remembered in
  \file{\textasciitilde/.config/medorama/home} (Windows: \code{\%APPDATA\%\textbackslash
  medorama\textbackslash home});
\item \file{\textasciitilde/.local/share/medorama} (Windows: \code{\%LOCALAPPDATA\%\textbackslash
  medorama}).
\end{enumerate}
A command other than \cmd{init} needs a catalog there. Without one, it stops with the message
\code{no catalog in DIR} and names the two ways out: create a catalog with \cmd{medorama init},
or name one with \opt{-{}-home}. \cmd{serve} prints the folder it uses.

\paragraph{Rights.} The commands act directly on the database, as the operating-system user who
can read the home folder, with the rights of an administrator and the PHI permission. This is
safe because whoever can run them under the catalog's account could read the database file
anyway. The audit log records them as \code{local:}\textit{USER}, where \textit{USER} is the
environment variable \code{USER} (Windows: \code{USERNAME}). The web page and the HTTP API, on
the other hand, act as signed-in catalog users (\cref{chap:users-privacy}). Of the commands that
list data, only \cmd{list -{}-phi} shows identifiers.

\paragraph{Output.} Results are printed on the standard output as text, one line per item; the
listing commands also take \opt{-{}-json}, which prints the same data as indented JSON, for
scripts. Messages go to the standard error. While an import runs in a terminal, a line on the
standard error shows its progress (rows read, the percentage done, rows per second, the time
left); it is not shown when the standard error is not a terminal.

\section{Exit status and messages}
\label{sec:cli-status}

\cmd{medorama} exits with status 0 on success, and 1 on an error, after printing one message on
the standard error that begins with \code{medorama:} (for example
\code{medorama: no user jdoe}). An error in the arguments (an unknown option, a missing required
argument, a value of the wrong type) is reported by the argument parser with the usage of the
subcommand, and the status is 2. \opt{-h} and \opt{-V} exit with status 0.

The VitalStream viewers exit with status 1 on an error (\code{vitalstream-rs: ...},
\code{vitalstream-cpp: ...}, \code{vitalstream-py: ...}) and 0 otherwise; an error in their
arguments gives status 2 (\cmd{vitalstream-rs}, \cmd{vitalstream-py}) or the nonzero status of
the C++ argument parser (\cmd{vitalstream-cpp}). \cmd{muse-py} and \cmd{dicom-py} exit with
status 0 on success, 1 on an error, and 2 on an error in the arguments.

\section{Environment variables}
\label{sec:cli-environment}

\begin{xltabular}{\linewidth}{@{}P{5.2cm}L@{}}
\caption{The environment variables that \cmd{medorama} reads.}
\label{tab:cli-environment}\\
\toprule
{\tabhead Variable} & {\tabhead Meaning} \\
\midrule
\endfirsthead
\toprule
{\tabhead Variable} & {\tabhead Meaning} \\
\midrule
\endhead
\bottomrule
\endfoot
\code{MEDORAMA\_HOME} & The catalog's home folder, when \opt{-{}-home} is not given
  (\cref{sec:cli-running}). \\
\code{MEDORAMA\_DATABASE\_URL} & A PostgreSQL database, without its password; overrides the
  setting \code{database\_url}. \\
\code{MEDORAMA\_DATABASE\_PASSWORD} & The PostgreSQL password; overrides
  \code{database\_password\_file} (and \opt{-{}-password-file} of \cmd{copy-database}). \\
\code{MEDORAMA\_KEY} & The database key, for a build with encryption; overrides
  \code{key\_file}. A key given to a build without encryption is an error. \\
\code{MEDORAMA\_SOURCE\_PASSWORD\_}\textit{NAME} & The password of the connected database
  \textit{NAME}, when no password file is given: \textit{NAME} in capitals, every character
  other than a letter or a digit written \code{\_} (the connection \code{omop-cdm}:
  \code{MEDORAMA\_SOURCE\_PASSWORD\_OMOP\_CDM}). \\
\code{MEDORAMA\_DEBUG\_SQL} & Set (to any value): print each statement that PostgreSQL refuses,
  never its values. \\
\code{USER}, \code{USERNAME} & The operating-system user named in the audit log
  (\code{local:}\textit{USER}). \\
\code{HOSTNAME}, \code{COMPUTERNAME} & The machine's name in the names of the workers that run
  jobs (else \code{localhost}). \\
\end{xltabular}

An action's program does not inherit the variables whose names begin with \code{MEDORAMA\_} or
\code{PGPASS}, nor any whose name contains \code{PASSWORD}, \code{SECRET}, \code{TOKEN},
\code{PASSPHRASE} or \code{CREDENTIAL}, ends with \code{\_KEY}, or is \code{DATABASE\_URL}; it is
given \code{MEDORAMA\_JOB}, the path of its job's description.

\section{The commands of Medorama}
\label{sec:cli-medorama}

\Cref{tab:cli-commands} lists the subcommands of \cmd{medorama}; the entries that follow are
in the same order.

\begin{xltabular}{\linewidth}{@{}P{5.0cm}L@{}}
\caption{The subcommands of \cmd{medorama}.}
\label{tab:cli-commands}\\
\toprule
{\tabhead Command} & {\tabhead Does} \\
\midrule
\endfirsthead
\toprule
{\tabhead Command} & {\tabhead Does} \\
\midrule
\endhead
\bottomrule
\endfoot
\multicolumn{2}{@{}l}{\textit{The catalog}} \\
\cmd{init} & Create a catalog: settings, database and the first admin \\
\cmd{use} & Make the catalog of \opt{-{}-home} the default of later commands \\
\cmd{serve} & Run the web interface, the HTTP API and the background workers \\
\cmd{copy-database} & Copy a SQLite catalog into an empty PostgreSQL database \\
\midrule
\multicolumn{2}{@{}l}{\textit{Bringing data in}} \\
\cmd{import-table} & Import a database: CSV, Parquet or spreadsheet files \\
\cmd{import-clif} & Import a CLIF folder, each table a database \\
\cmd{formats} & List the import formats; check a format file \\
\cmd{imports list}, \cmd{imports remove} & The imports (batches); remove one with what it added \\
\cmd{sources list}, \cmd{sources remove} & The sources; remove everything one contributed \\
\cmd{demo add}, \cmd{demo remove} & Load the demo; remove it \\
\cmd{connections add}, \cmd{list}, \cmd{check}, \cmd{sync}, \cmd{remove} & Connected databases \\
\midrule
\multicolumn{2}{@{}l}{\textit{Encounters and recordings}} \\
\cmd{list} & List encounters \\
\cmd{recordings} & List recordings \\
\cmd{remove} & Remove recordings from the catalog (the files are untouched) \\
\cmd{remove-encounters} & Remove encounters with their bed stays and items \\
\cmd{check} & Re-check the recordings: missing, changed, index folders gone \\
\midrule
\multicolumn{2}{@{}l}{\textit{Actions and jobs}} \\
\cmd{actions} & List the installed actions \\
\cmd{run} & Queue actions on recordings \\
\cmd{work} & Run the queued jobs here until none is left \\
\cmd{jobs} & List jobs \\
\cmd{cancel} & Cancel queued or running jobs \\
\midrule
\multicolumn{2}{@{}l}{\textit{Add-ons, users, audit}} \\
\cmd{addons list}, \cmd{install}, \cmd{remove} & The adapters of recording formats \\
\cmd{user add}, \cmd{list}, \cmd{set}, \cmd{remove} & The catalog's users \\
\cmd{audit} & Show the audit log; verify or seal it \\
\end{xltabular}

\subsection{The catalog}
\label{sec:cli-catalog}

\apientry{medorama init}{Create a catalog}
\begin{lstlisting}[style=signature]
medorama init --results DIR [--index-dir DIR] [--admin NAME] [--timezone TZ] [--port N]
              [--password-stdin] [--database-url URL [--database-password-file FILE]]
              [--no-demo]
\end{lstlisting}
Creates a catalog in the home folder (\cref{sec:cli-running}): the folder itself, readable by
its owner only, with \file{medorama.toml} (commented) and a folder \file{actions/}; the results
folder and the index folder; the database, in SQLite (\file{medorama.sqlite} in the home folder)
or in PostgreSQL; and the first user, an admin with the PHI permission. It asks twice for that
user's password (at least 10 characters; three attempts). The database is made first and the
settings file last, so that a failed \cmd{init} leaves no half-made catalog. \cmd{init} then
remembers the home folder, so that later commands use this catalog without \opt{-{}-home}, and,
when the installer made the demo (\cmd{tools/install-medorama.sh -{}-demo}), loads it as the
first admin's.

It refuses a home folder that holds a \file{medorama.toml} already, a SQLite file left there by
an earlier catalog, a PostgreSQL database that has users already, an unknown time zone, and a
database URL that is not a PostgreSQL URL or that holds a password. A SQLite database that cannot
be created (on a network file system, whose locking SQLite cannot rely on) is removed again, with
advice to choose a home folder on a local disk.
\begin{params}
\item[--results DIR] The results folder: action results and, by default, the indexes. Required;
  never a folder that holds recordings. A relative path is made absolute.
\item[--index-dir DIR] The index folder. Default: \file{DIR/index} of \opt{-{}-results}. A
  viewer whose \code{cache\_dir} is this folder opens the recordings the catalog indexed without
  indexing them again.
\item[--admin NAME] The first administrator's user name. Default: \code{admin}. With proxy
  sign-in, give the name that the proxy will send, such as \code{jdoe@hospital.org}
  (\cref{sec:sign-in-accounts}).
\item[--timezone TZ] The time zone of times written without an offset, and of the times shown
  (an IANA name). Default: \code{America/New\_York}.
\item[--port N] The port of the web interface and the API. Default: 8090.
\item[--password-stdin] Read the admin's password from the first line of the standard input
  instead of asking for it (scripts).
\item[--database-url URL] An empty PostgreSQL database instead of SQLite
  (\code{postgres://user@host/db?sslmode=require}), without its password.
\item[--database-password-file FILE] The file that holds that database's password; only with
  \opt{-{}-database-url}.
\item[--no-demo] Leave out the demo that the installer made; \cmd{medorama demo add} loads it
  later. Without the installer's \opt{-{}-demo} there is nothing to leave out.
\end{params}
\begin{lstlisting}[style=shell]
medorama init --results /data/medorama-results
medorama init --home ~/medorama --results /data/medorama-results --timezone Europe/Paris
medorama init --results /data/medorama-results \
    --database-url "postgres://medorama@db.example.org/medorama" \
    --database-password-file /etc/medorama/db-password
\end{lstlisting}

\apientry{medorama use}{Make a catalog the default}
\begin{lstlisting}[style=signature]
medorama --home DIR use
\end{lstlisting}
Remembers the catalog of \opt{-{}-home} (which must hold a \file{medorama.toml}) as the one that
later commands use without \opt{-{}-home}, in \file{\textasciitilde/.config/medorama/home}
(Windows: \code{\%APPDATA\%\textbackslash medorama\textbackslash home}). It has no options.
\begin{lstlisting}[style=shell]
medorama --home ~/medorama-test use
\end{lstlisting}

\apientry{medorama serve}{Run the web interface and the workers}
\begin{lstlisting}[style=signature]
medorama serve [--host HOST] [--port N] [--no-workers]
\end{lstlisting}
Serves the web page and the HTTP API (\cref{chap:http-api}), and runs the background work: the
workers that run queued jobs (the setting \code{workers}), the automatic indexing of new and
changed recordings (\code{auto\_index}), the re-check of every recording every
\code{check\_interval\_s} seconds, and the sealing of the audit log every minute. It prints the
address to open, the catalog, the database, the waveform viewer it will start (on first use, in
the \ui{VitalStream} tab) or link to, the results and index folders, and the number of workers. On an
address other than \code{127.0.0.1}, \code{localhost} or \code{::1} without
\code{secure\_cookies = true}, it warns that the page is reachable from other machines without
HTTPS (\cref{chap:deployment}). It stops on \key{Ctrl}+\key{C} or on \code{SIGTERM} (a service
manager's stop): the jobs it was running go back to the queue, and the viewer it started is
stopped. An action that could not be loaded is reported, and the server runs without it.
\begin{params}
\item[--host HOST] The address to listen on. Default: the setting \code{host}
  (\code{127.0.0.1}).
\item[--port N] The port. Default: the setting \code{port} (8090).
\item[--no-workers] Run no jobs (and so no automatic indexing): another process runs
  \cmd{medorama work}.
\end{params}
\begin{lstlisting}[style=shell]
medorama serve
medorama serve --port 8095 --no-workers
\end{lstlisting}

\apientry{medorama copy-database}{Copy a SQLite catalog into PostgreSQL}
\begin{lstlisting}[style=signature]
medorama copy-database --to URL [--password-file FILE] [--ca-file FILE]
\end{lstlisting}
Copies the catalog, every table and every row, from its SQLite database into an empty
PostgreSQL database (\cref{chap:deployment}), and records the copy in the new database's audit
log. It prints the number of rows of each table, then the lines to add to \file{medorama.toml}
(\code{database\_url}, and \code{database\_password\_file} when one was given) for the catalog
to use it. The SQLite file is left as it was: keep it as a backup, then delete it securely. A
catalog that is in PostgreSQL already is refused.
\begin{params}
\item[--to URL] The PostgreSQL database (\code{postgres://user@host/db?sslmode=require}),
  without its password. Required.
\item[--password-file FILE] The file that holds its password. Default: the environment variable
  \code{MEDORAMA\_DATABASE\_PASSWORD}.
\item[--ca-file FILE] Certificate authorities (PEM) trusted for the server's certificate,
  besides the system's.
\end{params}
\begin{lstlisting}[style=shell]
medorama copy-database --to "postgres://medorama@db.example.org/medorama" \
    --password-file /etc/medorama/db-password
\end{lstlisting}

\subsection{Bringing data in}
\label{sec:cli-data}

\apientry{medorama import-table}{Import a database}
\begin{lstlisting}[style=signature]
medorama import-table FILE... [--format FORMAT] [--mapping FILE] [--category CAT]
                      [--name NAME] [--source SOURCE] [--dry-run]
\end{lstlisting}
Imports a database: one or more CSV files, Parquet files or spreadsheets (\file{.xlsx},
\file{.xls}, \file{.ods}; the first sheet), each read row by row, as one batch. The rows may hold
patients, encounters, bed stays, measurements, notes and links to files, in any combination
(\cref{chap:importing}). Columns are recognized by their names, or mapped with
\opt{-{}-mapping}; a known export is recognized by its columns and converted. Relative file paths
in the table are read from the folder of the first file. Several files must have the same
columns and make one database, which then needs a name (\opt{-{}-name}). It prints what was
added: rows, patients, encounters, items (with those already there and those waiting for their
encounter), recordings, attached files and rows held back, and the columns kept as the
database's own.
\begin{params}
\item[FILE...] The files of the table. At least one.
\item[--format FORMAT] A known export to convert: \code{auto} (the default: recognized by its
  columns), \code{none}, or the key of a format as \cmd{medorama formats} lists it: built in
  \code{ge-muse} (GE MUSE ECG metadata, the patient by MRN) and \code{ge-muse-patient-id} (by a
  patient ID), or one described in a file of \code{formats\_dirs}. Not case-sensitive.
\item[--mapping FILE] A column mapping, in JSON or TOML: catalog field = header, in TOML at the
  top level or under \code{[columns]} (\cref{chap:importing}). With a mapping, the names are not
  guessed: a header is read as a catalog field only if the mapping names it for that field, or if
  it is that field's own name (ignoring case) and the mapping names no other header for that
  field.
\item[--category CAT] The category of the table's items when it has no category column
  (\code{lab}, \code{ecg}, \code{note}, \dots; for recordings \code{icu}, \code{holter},
  \code{eeg} or \code{sleep}). Default: none.
\item[--name NAME] The batch's name. Default: the file's name; required with several files.
\item[--source SOURCE] The source label. Default: the batch's name. A \code{source} column of
  the table overrides it row by row.
\item[--dry-run] Check and count only; store nothing.
\end{params}
\begin{lstlisting}[style=shell]
medorama import-table demo/databases/patients.csv --dry-run
medorama import-table demo/databases/labs.csv --source DEMO
medorama import-table labs-2023.csv labs-2024.csv --name labs --category lab
\end{lstlisting}

\apientry{medorama import-clif}{Import a CLIF folder}
\begin{lstlisting}[style=signature]
medorama import-clif DIR [--source SOURCE] [--dry-run]
\end{lstlisting}
Imports a folder of the Common Longitudinal ICU data Format (its tables as CSV or Parquet
files): each table a database (\code{clif\_patient}, \code{clif\_vitals}, \dots), in order
(patients, hospitalizations, ADT, then the rest); a table in several files is one database
(\cref{chap:clif}). A file that is not a CLIF table Medorama reads is named and skipped. For each
table it prints the rows read and what was added, with up to five warnings. Unlike the page and
the API, the command does not check the folder against \code{data\_roots}.
\begin{params}
\item[DIR] The CLIF folder.
\item[--source SOURCE] The source label. Default: \code{clif}.
\item[--dry-run] Check and count only; store nothing.
\end{params}
\begin{lstlisting}[style=shell]
medorama import-clif /data/clif-synthetic --dry-run
\end{lstlisting}

\apientry{medorama formats}{List the import formats}
\begin{lstlisting}[style=signature]
medorama formats [--check FILE]
\end{lstlisting}
Lists the import formats, one per line: its key, its category, its label and where it comes from
(\code{built in}, or the TOML file in \code{formats\_dirs} that describes it). A format file that
cannot be read is reported on the standard error (\code{not read: ...}).
\begin{params}
\item[--check FILE] Instead, check a format file before installing it: print its key, its label
  and its columns, or the error.
\end{params}
\begin{lstlisting}[style=shell]
medorama formats
medorama formats --check my-export.toml
\end{lstlisting}

\apientry{medorama imports list}{List the imports}
\begin{lstlisting}[style=signature]
medorama imports list [--json]
\end{lstlisting}
Lists the imports (batches), newest first: number, kind, state, time, name, source, and the
recordings, encounters and EHR items it holds.
\begin{params}
\item[--json] Print JSON.
\end{params}
\begin{lstlisting}[style=shell]
medorama imports list
\end{lstlisting}

\apientry{medorama imports remove}{Remove an import}
\begin{lstlisting}[style=signature]
medorama imports remove ID [--purge-results] [--dry-run]
\end{lstlisting}
Removes one import and everything it added: recordings, patients, encounters, bed stays and EHR
items, and its values. What another import also added stays, and is counted as kept. Files on
disk are never touched.
\begin{params}
\item[ID] The import's number, as \cmd{imports list} shows it.
\item[--purge-results] Also delete the catalog's results and indexes of the recordings removed
  (only inside \code{results\_dir} and \code{index\_dir}).
\item[--dry-run] Only say what would be removed.
\end{params}
\begin{lstlisting}[style=shell]
medorama imports remove 12 --dry-run
\end{lstlisting}

\apientry{medorama sources list}{List the sources}
\begin{lstlisting}[style=signature]
medorama sources list [--json]
\end{lstlisting}
Lists the source labels and what each contributes: recordings, patients, encounters, EHR items,
and the number of imports.
\begin{params}
\item[--json] Print JSON.
\end{params}
\begin{lstlisting}[style=shell]
medorama sources list
\end{lstlisting}

\apientry{medorama sources remove}{Remove a source}
\begin{lstlisting}[style=signature]
medorama sources remove SOURCE [--purge-results] [--dry-run]
\end{lstlisting}
Removes everything that the source contributed, in every import, as \cmd{imports remove} does
for one import.
\begin{params}
\item[SOURCE] The source label, as \cmd{sources list} shows it.
\item[--purge-results] Also delete the catalog's results and indexes of the recordings removed.
\item[--dry-run] Only say what would be removed.
\end{params}
\begin{lstlisting}[style=shell]
medorama sources remove "ICU pilot" --dry-run
\end{lstlisting}

\apientry{medorama demo add}{Load the demo}
\begin{lstlisting}[style=signature]
medorama demo add [DIR]
\end{lstlisting}
Loads the demo, invented patients for trying Medorama: its tables as the source \code{DEMO}, its
OMOP database as the connection \code{demo\_omop}, and the cohort \code{DEMO} that holds its
patients (\cref{chap:administration}). It prints what it loaded.
\begin{params}
\item[DIR] A folder made by \cmd{tools/make-demo-data.py} (\cref{sec:cli-tools}). Default: the
  one the installer made with \opt{-{}-demo}, named in the file \file{medorama.demo} next to the
  programs; without either, it is an error.
\end{params}
\begin{lstlisting}[style=shell]
medorama demo add
medorama demo add /data/demo
\end{lstlisting}

\apientry{medorama demo remove}{Remove the demo}
\begin{lstlisting}[style=signature]
medorama demo remove [--dry-run]
\end{lstlisting}
Removes the demo and everything it added, as deleting the cohort \code{DEMO} in the page does:
its patients, encounters, bed stays, items and recordings (with the catalog's indexes and results
of them), the connection \code{demo\_omop} and the cohort. The demo's files stay.
\begin{params}
\item[--dry-run] Only say what would be removed.
\end{params}
\begin{lstlisting}[style=shell]
medorama demo remove --dry-run
\end{lstlisting}

\apientry{medorama connections add}{Connect a database}
\begin{lstlisting}[style=signature]
medorama connections add NAME --kind omop|sql --url URL [--password-file FILE] [--ca-file FILE]
    [--schema S] [--vocabulary-schema S] [--timezone TZ] [--identifier mrn|patient_id|none]
    [--mrn-authority A] [--timeout-s N] [--max-rows N] [--views FILE]
\end{lstlisting}
Connects a database that is read where it is, beside the catalog (\cref{chap:connecting}). It is
checked first: each view must answer. It prints the views and the data's version. Its persons are
linked to the catalog's patients when a search finds them; \cmd{connections sync} links them all
at once.
\begin{params}
\item[NAME] The connection's name, as criteria write it (\code{omop.lactate}).
\item[--kind KIND] \code{omop}: an OMOP CDM 5.3 or 5.4 database, through built-in views;
  \code{sql}: any layout, through the views of \opt{-{}-views}. Required.
\item[--url URL] \code{postgres://user@host/db?sslmode=require}, without a password, or
  \code{sqlite:/path/to/file} (a relative SQLite path is made absolute). Required.
\item[--password-file FILE] The file that holds the password. Default: the environment variable
  \code{MEDORAMA\_SOURCE\_PASSWORD\_}\textit{NAME} (\cref{sec:cli-environment}).
\item[--ca-file FILE] Certificate authorities (PEM) trusted for the server's certificate,
  besides the system's.
\item[--schema S] The schema of the CDM's tables. Default: none (the database's default).
\item[--vocabulary-schema S] The schema of OMOP's vocabulary, when it is apart from the CDM's.
  Default: none.
\item[--timezone TZ] The time zone of the database's times. Default: the catalog's time zone,
  written into the connection when it is added, and fixed from then on (changing the catalog's
  \code{timezone} later does not change it).
\item[--identifier ID] What OMOP's \code{person\_source\_value} holds: \code{mrn} (the default),
  \code{patient\_id} or \code{none}; with \code{none}, each person becomes a patient of its own,
  with the patient ID \code{<connection>:<person key>}.
\item[--mrn-authority A] The authority that issued its MRNs. Default: none.
\item[--timeout-s N] Each query's time limit, in seconds (PostgreSQL). Default: 60.
\item[--max-rows N] The number of rows a criterion may read, and of persons a sync may read.
  Default: 1000000.
\item[--views FILE] Kind \code{sql}: a TOML file with the views \code{persons}, \code{visits},
  \code{entries} and \code{version}, each as SQL.
\end{params}
\begin{lstlisting}[style=shell]
medorama connections add omop --kind omop \
    --url "postgres://reader@omop.example.org/cdm?sslmode=require" \
    --password-file /etc/medorama/omop-password --schema cdm --vocabulary-schema vocab \
    --mrn-authority HOSP
medorama connections add demo2 --kind omop --url sqlite:/data/demo/omop/omop.sqlite
medorama connections add warehouse --kind sql --url sqlite:/data/wh.sqlite \
    --views warehouse-views.toml
\end{lstlisting}

\apientry{medorama connections list}{List the connected databases}
\begin{lstlisting}[style=signature]
medorama connections list [--json]
\end{lstlisting}
Lists the connected databases: name, kind, URL (without a password), database, and the number of
persons linked.
\begin{params}
\item[--json] Print JSON.
\end{params}
\begin{lstlisting}[style=shell]
medorama connections list
\end{lstlisting}

\apientry{medorama connections check}{Check a connected database}
\begin{lstlisting}[style=signature]
medorama connections check NAME
\end{lstlisting}
Checks that the connected database answers: the time it took, each view, and its version.
\begin{lstlisting}[style=shell]
medorama connections check omop
\end{lstlisting}

\apientry{medorama connections sync}{Link every person}
\begin{lstlisting}[style=signature]
medorama connections sync NAME
\end{lstlisting}
Links every person of the connected database to the catalog's patients, by the rules of an
imported table. It reads every person in one query, within the connection's \code{max\_rows}: a
database with more persons than \code{max\_rows} fails to sync (raise \code{max\_rows}, or let
searches link the persons they find). It prints the number of persons, and of those linked; the
others are held for review, or have no valid identifier.
\begin{lstlisting}[style=shell]
medorama connections sync omop
\end{lstlisting}

\apientry{medorama connections remove}{Disconnect a database}
\begin{lstlisting}[style=signature]
medorama connections remove NAME
\end{lstlisting}
Disconnects the database: its links and what it added to the catalog are removed. The database
itself is not touched.
\begin{lstlisting}[style=shell]
medorama connections remove omop
\end{lstlisting}

\subsection{Encounters and recordings}
\label{sec:cli-encounters}

The options \opt{-{}-q} and \opt{-{}-search} take a search in the syntax of the page's search
boxes (\cref{chap:patients}): words, \code{field:value} terms (\code{unit:ICU7},
\code{admit:2023-12}), quotes around values with spaces, every term to match.

\apientry{medorama list}{List encounters}
\begin{lstlisting}[style=signature]
medorama list [--q SEARCH] [--unit UNIT] [--status STATUS] [--limit N] [--phi] [--json]
\end{lstlisting}
Lists the encounters that match: the total, then one line each with the status of its recordings,
the encounter and patient IDs, admission and discharge (local time), the units and the number of
recordings.
\begin{params}
\item[--q SEARCH] Search words (encounter ID, patient ID; with \opt{-{}-phi}, also MRN, name,
  date of birth). Default: every encounter.
\item[--unit UNIT] Only the encounters in this unit.
\item[--status STATUS] Only the encounters whose recordings are \code{ready}, \code{new},
  \code{queued}, \code{indexing}, \code{stale}, \code{failed}, \code{missing}, or which have
  \code{none}.
\item[--limit N] At most \textit{N} encounters. Default: 100.
\item[--phi] Show identifiers (MRN, name, date of birth, encounter number), and let the search
  match them.
\item[--json] Print JSON: \code{\{"total": ..., "items": [...]\}}.
\end{params}
\begin{lstlisting}[style=shell]
medorama list --unit ICU7 --status ready
medorama list --q P900001 --json
\end{lstlisting}

\apientry{medorama recordings}{List recordings}
\begin{lstlisting}[style=signature]
medorama recordings [--q SEARCH] [--unlinked] [--limit N] [--json]
\end{lstlisting}
Lists the recordings that match: the total, then one line each with its number, status (with the
percentage done while it is indexed), kind, format, start time and path, and its error if it has
one.
\begin{params}
\item[--q SEARCH] A search (fields \code{file}, \code{path}, \code{kind}, \code{format},
  \code{unit}, \code{bed}, \code{start}, \code{status}, \code{source}, \code{encounter}). Default:
  every recording.
\item[--unlinked] Only the recordings not linked to an encounter.
\item[--limit N] At most \textit{N} recordings. Default: 100.
\item[--json] Print JSON: \code{\{"total": ..., "items": [...]\}}.
\end{params}
\begin{lstlisting}[style=shell]
medorama recordings --unlinked
medorama recordings --q "status:failed"
\end{lstlisting}

\apientry{medorama remove}{Remove recordings from the catalog}
\begin{lstlisting}[style=signature]
medorama remove [ID|PATH...] [--search SEARCH] [--purge-results] [--dry-run]
\end{lstlisting}
Removes recordings from the catalog: the catalog's entries only. The files are never touched.
A path that is not in the catalog is an error.
\begin{params}
\item[ID or PATH] Recordings by number, or by path.
\item[--search SEARCH] Also the recordings that match a search.
\item[--purge-results] Also delete the catalog's results and indexes of them (only inside
  \code{results\_dir} and \code{index\_dir}).
\item[--dry-run] Only say how many would be removed.
\end{params}
\begin{lstlisting}[style=shell]
medorama remove 41 42 --dry-run
medorama remove --search "unit:ICU7 start:2023-12" --purge-results
\end{lstlisting}

\apientry{medorama remove-encounters}{Remove encounters}
\begin{lstlisting}[style=signature]
medorama remove-encounters [ENCOUNTER...] [--search SEARCH] [--all] [--dry-run]
\end{lstlisting}
Removes encounters with their bed stays and EHR items. Their recordings stay in the catalog,
unlinked.
\begin{params}
\item[ENCOUNTER] Encounters by encounter ID. An unknown ID is an error.
\item[--search SEARCH] Also the encounters that match a search.
\item[--all] Every encounter.
\item[--dry-run] Only say what would be removed.
\end{params}
\begin{lstlisting}[style=shell]
medorama remove-encounters E9100002 --dry-run
\end{lstlisting}

\apientry{medorama check}{Re-check the recordings}
\begin{lstlisting}[style=signature]
medorama check
\end{lstlisting}
Re-checks every recording, as the server does every \code{check\_interval\_s}: whether its files
are missing, are back, or changed, and whether its index folder is still there. It prints the
counts. It has no options.
\begin{lstlisting}[style=shell]
medorama check
\end{lstlisting}

\subsection{Actions and jobs}
\label{sec:cli-actions}

An action is a program run on recordings: the built-in \code{index} and \code{identify}, and the
plug-ins of \code{actions\_dirs} (\cref{chap:actions}). \cmd{run} queues jobs; \cmd{serve}'s
workers, or \cmd{work}, run them.

\apientry{medorama actions}{List the installed actions}
\begin{lstlisting}[style=signature]
medorama actions
\end{lstlisting}
Lists the actions: name, version, title, the actions each requires, and why one cannot run; then
the plug-ins that could not be loaded. It has no options.
\begin{lstlisting}[style=shell]
medorama actions
\end{lstlisting}

\apientry{medorama run}{Queue actions}
\begin{lstlisting}[style=signature]
medorama run ACTION... (--all | --encounter ID... | --recording ID... | --search SEARCH |
             --cohort N) [--force] [--param ACTION.KEY=VALUE]... [--priority N]
\end{lstlisting}
Queues jobs of the actions on the recordings chosen, and prints the run's number and how many
jobs were queued, already done, already queued and not applicable. At least one of
\opt{-{}-all}, \opt{-{}-encounter}, \opt{-{}-recording}, \opt{-{}-search} and \opt{-{}-cohort} is
needed.
\begin{params}
\item[ACTION...] The actions' names (\code{index}, \dots). At least one.
\item[--all] Every recording.
\item[--encounter ID] The recordings of an encounter, by encounter ID; repeatable.
\item[--recording ID] A recording, by number; repeatable.
\item[--search SEARCH] The recordings of the encounters that match a search.
\item[--cohort N] The recordings of the encounters of a cohort's members, by the cohort's
  number as the \ui{Cohorts} tab shows it.
\item[--force] Run again where an action was done already.
\item[--param ACTION.KEY=VALUE] A parameter of an action; \textit{VALUE} is read as JSON, else
  as text. Repeatable.
\item[--priority N] The jobs' priority: higher runs first. Default: 0.
\end{params}
\begin{lstlisting}[style=shell]
medorama run index --all
medorama run index --encounter E9100002 --force
medorama run recording-summary --cohort 3 --param recording-summary.max_channels=16 --priority 5
\end{lstlisting}

\apientry{medorama work}{Run the queue here}
\begin{lstlisting}[style=signature]
medorama work [--workers N]
\end{lstlisting}
Runs queued jobs in this process until none is left, then prints how many succeeded and failed
here. Several \cmd{work} processes, on one or several machines, share the queue of a PostgreSQL
catalog.
\begin{params}
\item[--workers N] Jobs at the same time. Default: the setting \code{workers}; at least 1.
\end{params}
\begin{lstlisting}[style=shell]
medorama work --workers 4
\end{lstlisting}

\apientry{medorama jobs}{List jobs}
\begin{lstlisting}[style=signature]
medorama jobs [--state STATE] [--limit N] [--json]
\end{lstlisting}
Lists jobs, running ones first, then queued ones (higher priority first), then the others, newest
first: number, state, action, progress, recording and its error.
\begin{params}
\item[--state STATE] Only the jobs in this state: \code{queued}, \code{running}, \code{done},
  \code{failed} or \code{cancelled}.
\item[--limit N] At most \textit{N} jobs. Default: 50.
\item[--json] Print JSON.
\end{params}
\begin{lstlisting}[style=shell]
medorama jobs --state failed
\end{lstlisting}

\apientry{medorama cancel}{Cancel jobs}
\begin{lstlisting}[style=signature]
medorama cancel (--job N | --run N | --all)
\end{lstlisting}
Cancels queued and running jobs, and prints how many. One of the options is needed; \opt{-{}-job}
and \opt{-{}-run} together cancel that job only if it belongs to that run.
\begin{params}
\item[--job N] One job, by number.
\item[--run N] The unfinished jobs of a run.
\item[--all] Every unfinished job.
\end{params}
\begin{lstlisting}[style=shell]
medorama cancel --run 7
\end{lstlisting}

\subsection{Add-ons}
\label{sec:cli-addons}

Add-ons are the adapters that read the recording formats Medorama does not read itself
(\cref{chap:addons}). They are installed in the adapter folder, the setting \code{adapters\_dir}.
Installing is done on the server, from the command line only, never through the page; installs
and removals are written to the audit log. A server that runs already sees the change after a
restart.

\apientry{medorama addons list}{List the installed adapters}
\begin{lstlisting}[style=signature]
medorama addons list [--json]
\end{lstlisting}
Lists the adapters of the adapter folder: name, version, implementation, the formats each reads
(label and file extensions) and its program.
\begin{params}
\item[--json] Print JSON.
\end{params}
\begin{lstlisting}[style=shell]
medorama addons list
\end{lstlisting}

\apientry{medorama addons install}{Install an adapter}
\begin{lstlisting}[style=signature]
medorama addons install PROGRAM
\end{lstlisting}
Copies an adapter's program into the adapter folder (made executable), after checking that it is
named \code{<name>-adapter} (\code{<name>-adapter.exe} on Windows) and that it describes itself
in this protocol. An adapter of the same name is replaced.
\begin{params}
\item[PROGRAM] The adapter's program.
\end{params}
\begin{lstlisting}[style=shell]
medorama addons install example-adapter-0.1.0-linux-x86_64/example-adapter
\end{lstlisting}

\apientry{medorama addons remove}{Remove an adapter}
\begin{lstlisting}[style=signature]
medorama addons remove NAME
\end{lstlisting}
Removes an installed adapter, by its name as \cmd{addons list} shows it.
\begin{lstlisting}[style=shell]
medorama addons remove example
\end{lstlisting}

\subsection{Users and the audit log}
\label{sec:cli-users}

Users, roles and the PHI permission are described in \cref{chap:users-privacy}, and the
sign-in methods in \cref{chap:sign-in}. A user name has 1 to 64 characters: letters, digits, and
\code{.}, \code{\_}, \code{-} and \code{@}. A password has at least 10 characters. The command
asks for it twice (with three attempts), or reads it from the first line of the standard input
with \opt{-{}-password-stdin}. Every change is written to the audit log.

\apientry{medorama user add}{Add a user}
\begin{lstlisting}[style=signature]
medorama user add USERNAME --role viewer|analyst|admin [--phi] [--display-name NAME]
                  [--password-stdin | --no-password]
\end{lstlisting}
Adds a user. Without \opt{-{}-password-stdin} or \opt{-{}-no-password}, the command asks for the
password.
\begin{params}
\item[USERNAME] The user name, for signing in. With proxy sign-in, it is the exact name that the
  institution's proxy sends, such as \code{jdoe@hospital.org} (case does not matter).
\item[--role ROLE] \code{viewer}, \code{analyst} or \code{admin}. Required.
\item[--phi] The PHI permission: the user may see and search identifiers.
\item[--display-name NAME] The name shown. Default: none.
\item[--password-stdin] Read the password from the standard input.
\item[--no-password] Give the account no password of its own: it gets a random password that
  nobody knows. Use it for accounts that sign in through the institution's sign-in
  (\code{sign\_in = "proxy"}; \cref{sec:sign-in-accounts}). With password sign-in, such an
  account cannot sign in until an administrator sets a password. It cannot be combined with
  \opt{-{}-password-stdin}.
\end{params}
\begin{lstlisting}[style=shell]
medorama user add analyst1 --role analyst --display-name "Analyst One"
medorama user add jdoe@hospital.org --role viewer --no-password   # proxy sign-in
\end{lstlisting}

\apientry{medorama user list}{List users}
\begin{lstlisting}[style=signature]
medorama user list
\end{lstlisting}
Lists the users: name, role, \code{PHI} for the PHI permission, \code{disabled}, and the display
name.
\begin{lstlisting}[style=shell]
medorama user list
\end{lstlisting}

\apientry{medorama user set}{Change a user}
\begin{lstlisting}[style=signature]
medorama user set USERNAME [--role ROLE] [--phi true|false] [--disabled true|false]
                  [--display-name NAME] [--password | --password-stdin]
\end{lstlisting}
Changes what the options give, and leaves the rest as it is.
\begin{params}
\item[--role ROLE] \code{viewer}, \code{analyst} or \code{admin}.
\item[--phi BOOL] Give (\code{true}) or take away (\code{false}) the PHI permission.
\item[--disabled BOOL] Disable (\code{true}) or enable (\code{false}) the account.
\item[--display-name NAME] The name shown.
\item[--password] Set a new password, asked for.
\item[--password-stdin] Set a new password, read from the standard input.
\end{params}
\begin{lstlisting}[style=shell]
medorama user set analyst1 --phi true
medorama user set analyst1 --disabled true
\end{lstlisting}

\apientry{medorama user remove}{Delete a user}
\begin{lstlisting}[style=signature]
medorama user remove USERNAME
\end{lstlisting}
Deletes a user.
\begin{lstlisting}[style=shell]
medorama user remove analyst1
\end{lstlisting}

\apientry{medorama audit}{Show, verify or seal the audit log}
\begin{lstlisting}[style=signature]
medorama audit [--limit N]
medorama audit --verify
medorama audit --seal
medorama audit --seal --verify
\end{lstlisting}
Shows the newest entries of the audit log: time, user, event, target and detail
(\cref{chap:users-privacy}). With \opt{-{}-verify} or \opt{-{}-seal} it does that instead; with
both, it seals, then verifies every seal, the new one too.
\begin{params}
\item[--limit N] The number of entries shown. Default: 50.
\item[--verify] Check the seals: whether a sealed entry was changed, deleted or added. It prints
  the seals, the entries sealed and those not sealed yet; a broken chain is an error (status 1).
\item[--seal] Seal the entries not sealed yet, now (the server does it every minute), and append
  them to \code{audit\_export} when it is set.
\end{params}
\begin{lstlisting}[style=shell]
medorama audit --limit 20
medorama audit --verify
medorama audit --seal --verify
\end{lstlisting}

\section{The VitalStream viewers}
\label{sec:cli-vitalstream}

The three VitalStream viewers, in Rust, C++ and Python, have the same options, the same settings
file (\cref{chap:settings-reference}), the same API (\cref{chap:viewer-api}) and the same password
file; they differ in their default ports and in a few details
(\cref{tab:cli-viewer-differences}).

\apientry{vitalstream-rs, vitalstream-cpp, vitalstream-py}{Run a VitalStream viewer}
\begin{lstlisting}[style=signature]
vitalstream-rs [--config FILE] [--host HOST] [--port N] [--timezone TZ] [--root DIR]...
               [--cache-dir DIR] [--adapters-dir DIR] [--qc-verbose] [--no-browser]
               [--no-auth] [--set-password | --clear-password | --index PATH] [-V]
\end{lstlisting}
Starts the viewer's web server and opens its page in a browser, after a second. It prints the
settings file it read, its address (an IPv6 host in brackets: \code{http://[::1]:8050/}), its
root, its data roots, the index folder, the time zone, and whether a password protects it; a data root outside the viewer's root is left out with a warning. With
\opt{-{}-set-password}, \opt{-{}-clear-password} or \opt{-{}-index} it does that and exits,
without serving. The options override the settings file, which overrides the defaults.
\begin{params}
\item[--config FILE] The settings file (TOML). Default: the first that exists of
  \file{viewers/vitalstream/vitalstream.toml} (only for a viewer run from the build folder of a
  clone) and \file{\textasciitilde/.config/vitalstream/vitalstream.toml}; a
  \file{vitalstream.toml} in the working folder is never read unless \opt{-{}-config} names it.
  The viewer says on start which file it read. Its top-level keys apply to the
  three viewers, those of its section \code{[rs]}, \code{[cpp]} or \code{[py]} to that viewer
  only.
\item[--host HOST] The address to listen on. Default: \code{127.0.0.1} (this machine only).
\item[--port N] The port. Default: 8060 (\cmd{vitalstream-rs}), 8070 (\cmd{vitalstream-cpp}),
  8050 (\cmd{vitalstream-py}).
\item[--timezone TZ] The display time zone, an IANA name. Default: \code{America/New\_York}.
\item[--root DIR] A folder the viewer may open, inside the viewer's root (fixed when it was
  built); repeatable, and in place of the setting \code{data\_roots}. Default: the viewer's root.
\item[--cache-dir DIR] The index folder. Default: \file{sample\_data/outputs/index} of the
  viewer's folder when it runs from a clone of the repository, else
  \file{\textasciitilde/.cache/vitalstream}.
\item[--adapters-dir DIR] The folder of adapters. Default:
  \file{\textasciitilde/.local/share/medorama/adapters}.
\item[--qc-verbose] Also mark the irregularities a reader repaired (data blocks put back in
  order).
\item[--no-browser] Do not open a browser window (on a server).
\item[--no-auth] Run without asking for the password, even when one is set.
\item[--set-password] Set (or change) the password that the page asks for, at least 8
  characters, asked twice; then exit. The viewer must be restarted for it to take effect. A
  viewer built without password support refuses (\code{this viewer was built without password
  support}); it never asks for a password, and ignores one that is set already.
\item[--clear-password] Remove the password, then exit. A viewer built without password support
  refuses.
\item[--index PATH] Index the recording \textit{PATH} (a WFDB record, collection or folder of
  records, an EDF, EDF+, BDF or BDF+ file, or a file that an installed adapter reads) into the
  index folder, print
  \code{indexed PATH -> CACHE/KEY/}, and exit.
\item[-V, --version] Print the program's name and version (\code{vitalstream-py 0.1.0}), then
  exit.
\end{params}
\begin{lstlisting}[style=shell]
vitalstream-cpp
vitalstream-cpp --set-password
vitalstream-rs --no-browser --port 8061 --root ~/icu
vitalstream-py --index ~/icu/demo-record
\end{lstlisting}

A viewer whose settings say \code{require\_password = true} refuses to start without a password
(Medorama sets it for the viewer it runs). The password is kept in \code{auth\_file}, by default
\file{\textasciitilde/.config/vitalstream/auth.json}, shared by the three viewers and by
\cmd{muse-py serve} and \cmd{dicom-py serve}.

\begin{xltabular}{\linewidth}{@{}P{4.2cm}LLL@{}}
\caption{Where the three VitalStream viewers' command lines differ.}
\label{tab:cli-viewer-differences}\\
\toprule
{\tabhead } & {\tabhead vitalstream-rs} & {\tabhead vitalstream-cpp} & {\tabhead vitalstream-py} \\
\midrule
\endfirsthead
\toprule
{\tabhead } & {\tabhead vitalstream-rs} & {\tabhead vitalstream-cpp} & {\tabhead vitalstream-py} \\
\midrule
\endhead
\bottomrule
\endfoot
Default port & 8060 & 8070 & 8050 \\
Section of the settings file & \code{[rs]} & \code{[cpp]} & \code{[py]} \\
Argument errors & status 2 & the C++ parser's status & status 2 \\
\end{xltabular}

\section{The ECG and PACS database tools}
\label{sec:cli-ecg-pacs}

\cmd{muse-py} reads GE MUSE resting ECG XML files; \cmd{dicom-py} reads DICOM files. Each builds
databases to import into Medorama (\cref{chap:ecg-imaging}) and converts ECGs to WFDB records, and
each \cmd{serve} runs a viewer from this computer: the ECG viewer (\cmd{muse-py}) and the PACS
viewer (\cmd{dicom-py}). Both take \opt{-h} and \opt{-{}-version}.

\begin{phibox}
The databases that \cmd{muse-py index} and \cmd{dicom-py index} write hold the patient's MRN or
ID\phimark{}, the key that links each row to its patient: keep them as protected health
information. The JSON and WFDB outputs carry no patient identity.
\end{phibox}

\apientry{muse-py json}{An ECG as the ECG viewer's JSON}
\begin{lstlisting}[style=signature]
muse-py json FILE
\end{lstlisting}
Prints the ECG of a MUSE XML file in the ECG viewer's JSON form, on the standard output.
\begin{lstlisting}[style=shell]
muse-py json ecg/afib_muse.xml > afib.json
\end{lstlisting}

\apientry{muse-py wfdb}{An ECG as a WFDB record}
\begin{lstlisting}[style=signature]
muse-py wfdb FILE OUT_DIR [--name NAME]
\end{lstlisting}
Writes the ECG as a WFDB record in \textit{OUT\_DIR} (created if needed): a signal file in format
16, 1000 units per millivolt, and a header whose comments carry the details, the measurements and
the statements. It prints the header's path.
\begin{params}
\item[--name NAME] The record's name. Default: the file's name without its extension.
\end{params}
\begin{lstlisting}[style=shell]
muse-py wfdb ecg/afib_muse.xml records/ --name afib
\end{lstlisting}

\apientry{muse-py index}{A database of a folder of MUSE files}
\begin{lstlisting}[style=signature]
muse-py index FOLDER -o FILE
\end{lstlisting}
Writes an ECG database to import: one row per MUSE file under \textit{FOLDER} (searched in every
subfolder), with the MUSE PatientID in a column named \code{mrn} (Medorama reads it as an MRN),
the time, the category \code{ecg}, the ventricular rate, every measurement as a column, the
computer's and the confirmed interpretation, and the file's path, relative to the folder of the
CSV. It prints the number of ECGs, and names the files it could not
read on the standard error.
\begin{params}
\item[-o, --output FILE] The CSV file to write. Required.
\end{params}
\begin{lstlisting}[style=shell]
muse-py index /data/demo/ecg -o /data/demo/ecg-muse.csv
\end{lstlisting}

\apientry{dicom-py index}{A database of a folder of DICOM files}
\begin{lstlisting}[style=signature]
dicom-py index FOLDER -o FILE
\end{lstlisting}
Writes a database to import: one row per loop, run, series, report, PDF or ECG under
\textit{FOLDER}, with the patient ID, the time, the category, the measurements and the file,
relative to the folder of the CSV. It prints the number of objects, and names those it could not
link on the standard error.
\begin{params}
\item[-o, --output FILE] The CSV file to write. Required.
\end{params}
\begin{lstlisting}[style=shell]
dicom-py index /data/demo/echo -o /data/demo/echo.csv
\end{lstlisting}

\apientry{dicom-py wfdb}{A DICOM ECG as a WFDB record}
\begin{lstlisting}[style=signature]
dicom-py wfdb FILE OUT_DIR [--name NAME]
\end{lstlisting}
Writes a DICOM ECG (12-lead or Holter) as a WFDB record in \textit{OUT\_DIR}, and prints its
header's path.
\begin{params}
\item[--name NAME] The record's name. Default: the file's name.
\end{params}
\begin{lstlisting}[style=shell]
dicom-py wfdb ecg/ecg_12lead.dcm records/
\end{lstlisting}

\apientry{muse-py serve, dicom-py serve}{Serve the ECG or PACS viewer}
\begin{lstlisting}[style=signature]
muse-py serve [--root DIR] [--host HOST] [--port N] [--auth-file FILE] [--no-auth]
              [--no-browser] [--set-password | --clear-password]
dicom-py serve (the same options)
\end{lstlisting}
Serves the viewer's page from this computer, with the files and folders under a root, and opens
\code{http://HOST:PORT/browse} in a browser after a second (the server's folders). Only what lies
under the root, every symbolic link resolved, is listed and served, and never a hidden entry (a
name that begins with \code{.}) (\cref{chap:viewer-api}). It
prints the address, the root, and whether a password protects it; \key{Ctrl}+\key{C} stops it.
\begin{params}
\item[--root DIR] The folder served. Default: the home folder. It must be a folder inside the
  root the package was installed with (\code{VIEWER\_ROOT}; by default the whole file system): a
  \opt{-{}-root} that contains that root is narrowed to it, and one outside it is refused (status 1).
\item[--host HOST] The address to listen on. Default: \code{127.0.0.1}. A request must name the
  server by a loopback name, an IP address or this host; any other name is refused (no DNS
  rebinding).
\item[--port N] The port. Default: 8080 (\cmd{muse-py}), 8081 (\cmd{dicom-py}).
\item[--auth-file FILE] The password file. Default:
  \file{\textasciitilde/.config/vitalstream/auth.json}, shared with VitalStream.
\item[--no-auth] Run without asking for the password.
\item[--no-browser] Do not open a browser window.
\item[--set-password] Set (or change) the password that the page asks for, at least 8
  characters, asked twice; then exit. For \cmd{muse-py serve} and \cmd{dicom-py serve}, a
  sign-in lasts 30 days, until the server stops, or until a browser signs out, which ends every
  sign-in; after a wrong password the next try, from any browser, waits 1\,s, then 2, 4 \dots{}
  up to 60\,s. A package installed with
  \code{VIEWER\_PASSWORD=0} refuses (status 1); it never asks for a password, and ignores one
  that is set already.
\item[--clear-password] Remove the password file, then exit. A package installed with
  \code{VIEWER\_PASSWORD=0} refuses (status 1).
\end{params}
\begin{lstlisting}[style=shell]
muse-py serve
dicom-py serve --root /data/demo/echo --no-browser
ssh -N -L 8081:127.0.0.1:8081 user@server        # then, on your computer, open the address
\end{lstlisting}

\section{The repository's tools}
\label{sec:cli-tools}

The folder \file{tools/} of the repository holds the installers, generators and checks; none of
it is part of a package. They are run from a clone (\cref{chap:installation}).

\apientry{tools/install-medorama.sh}{Install Medorama and one viewer}
\begin{lstlisting}[style=signature]
tools/install-medorama.sh [--viewer cpp|rs|py|none] [--demo] [--no-rustup] [--no-tools]
powershell -ExecutionPolicy Bypass -File tools\install-medorama.ps1 [-Viewer cpp|rs|py|none]
    [-Demo] [-NoRustup] [-NoTools]
\end{lstlisting}
Builds Medorama and one VitalStream viewer, without administrator rights (macOS and Linux; the
\file{.ps1} on Windows), and installs the database tools \cmd{muse-py} and \cmd{dicom-py}. The
programs go into \file{.venv/bin} (Windows \file{.venv\textbackslash Scripts}), and the viewer
chosen is recorded in \file{medorama.viewer} there, so that Medorama starts that viewer. An unknown
option is an error (status 2).
\begin{params}
\item[--viewer VIEWER] \code{cpp} (the default: the C++ viewer), \code{rs}, \code{py}, or
  \code{none} (Medorama alone).
\item[--demo] Also make the demo, in \file{.venv/share/medorama/demo}, named in
  \file{medorama.demo} next to the programs: \cmd{medorama init} then loads it. Off by default.
\item[--no-rustup] Fail rather than install Rust when \cmd{cargo} is missing.
\item[--no-tools] Leave out \cmd{muse-py} and \cmd{dicom-py}.
\item[-h, --help] Print the usage.
\end{params}
\begin{lstlisting}[style=shell]
tools/install-medorama.sh
tools/install-medorama.sh --viewer rs --demo
\end{lstlisting}

\apientry{tools/install.sh}{Install everything, for development}
\begin{lstlisting}[style=signature]
tools/install.sh [--py] [--rs] [--cpp] [--demo] [--no-rustup]
powershell -ExecutionPolicy Bypass -File tools\install.ps1 [-Py] [-Rs] [-Cpp] [-Demo] [-NoRustup]
\end{lstlisting}
Installs into the clone, for development: \file{.venv} with the Python packages (the Python
viewer, \cmd{muse-py}, \cmd{dicom-py} and the test tools, editable), the Rust workspace
(\cmd{medorama}, \cmd{vitalstream-rs}) and the C++ project (\cmd{vitalstream-cpp} and its tests),
and enables the repository's git hooks. Without \opt{-{}-py}, \opt{-{}-rs} or \opt{-{}-cpp}: all
of it. When \cmd{cargo} is missing, both installers install Rust with rustup (on Windows,
\file{rustup-init.exe} is checked against its published SHA-256 first).
\begin{params}
\item[--py] The Python packages only.
\item[--rs] The Python packages and the Rust workspace.
\item[--cpp] The Python packages and the C++ viewer; \opt{-{}-rs -{}-cpp}: both.
\item[--demo] Also make the demo, as \cmd{install-medorama.sh -{}-demo} does.
\item[--no-rustup] Fail rather than install Rust when \cmd{cargo} is missing.
\item[-h, --help] Print the usage.
\end{params}
\begin{lstlisting}[style=shell]
tools/install.sh --rs
\end{lstlisting}

\apientry{tools/install-gcc.sh}{A C++20 compiler without administrator rights}
\begin{lstlisting}[style=signature]
tools/install-gcc.sh [--prefix DIR]
\end{lstlisting}
On Linux, installs GCC 13, CMake, Ninja and Python 3.12 from conda-forge, in a Miniforge
installation (environment \code{vs-build}), for machines whose compiler is too old for the C++
viewer. The Miniforge installer is checked against its published SHA-256 before it runs.
Afterwards, \cmd{source \textasciitilde/miniforge3/bin/activate vs-build} in each shell
that builds it.
\begin{params}
\item[--prefix DIR] The Miniforge folder. Default: \file{\textasciitilde/miniforge3}; a Miniforge
  or Anaconda installation there is reused.
\end{params}
\begin{lstlisting}[style=shell]
tools/install-gcc.sh
\end{lstlisting}

\apientry{tools/fetch-cpp-deps.sh}{Download the C++ dependencies}
\begin{lstlisting}[style=signature]
tools/fetch-cpp-deps.sh
powershell -ExecutionPolicy Bypass -File tools\fetch-cpp-deps.ps1
\end{lstlisting}
Downloads the C++ dependencies pinned in \file{viewers/vitalstream/cpp/cmake/deps.cmake} into
\file{viewers/vitalstream/cpp/deps/}, checking each archive's SHA-256, so that a build needs no
internet access. It has no options.
\begin{lstlisting}[style=shell]
tools/fetch-cpp-deps.sh
\end{lstlisting}

\apientry{tools/gen\_timezones.py}{The C++ viewer's time zones}
\begin{lstlisting}[style=signature]
python3 tools/gen_timezones.py
\end{lstlisting}
Generates \file{viewers/vitalstream/cpp/vitalstream/src/timezones.cpp}, the IANA time-zone names
that \cmd{vitalstream-cpp} accepts, from this Python's \code{zoneinfo} database. It has no
options.

\apientry{tools/make-demo-data.py}{Make the synthetic demo}
\begin{lstlisting}[style=signature]
.venv/bin/python tools/make-demo-data.py [OUT] [--examples DIR] [--dicom DIR]
\end{lstlisting}
Writes a synthetic data set under \textit{OUT} (default: \file{sample\_data/demo} of the clone):
ICU recordings as WFDB collections, a 24-hour Holter, an overnight sleep study (an EDF+ file,
\file{sleep/P900004\_psg\_20240305.edf}, linked by \file{sleep\_index.csv}), 12-lead ECGs, DICOM
imaging, echo and angiography, the databases that link them to invented patients, and an OMOP
database to connect. Everything is
invented. \cmd{medorama demo add OUT} loads it. The same options give the same files, byte for
byte, on every run. It needs NumPy and wfdb-x.
\begin{params}
\item[OUT] The folder to write.
\item[--examples DIR] Also give local, already de-identified ECG examples to some demo patients;
  MUSE XML files are de-identified again on the way.
\item[--dicom DIR] Also give local public DICOM sample studies to the demo patients they fit, in
  place of the synthetic CT series, de-identified on the way.
\item[-h, --help] Print the usage.
\end{params}
\begin{lstlisting}[style=shell]
.venv/bin/python tools/make-demo-data.py /data/demo
\end{lstlisting}

\apientry{tools/make-viewer-examples.py}{The viewers' example files}
\begin{lstlisting}[style=signature]
python3 tools/make-viewer-examples.py
\end{lstlisting}
Writes the example files of the ECG and PACS viewers (\file{viewers/ecg/examples/},
\file{viewers/pacs/examples/}) with the demo's generators: a 12-lead WFDB record, a MUSE XML ECG,
a DICOM ECG, a PDF and a PNG printout; an echo loop, a coronary angiography run and a chest X-ray.
Every value is invented. It needs NumPy and wfdb-x, and has no options.

\apientry{tools/screenshots.py}{The documentation's screenshots}
\begin{lstlisting}[style=signature]
python3 tools/screenshots.py [OUT] [--url URL] [--user NAME] --password-file FILE --demo DIR
\end{lstlisting}
Takes the documentation's screenshots with headless Chrome or Chromium (the environment variable
\code{MEDORAMA\_TEST\_CHROME}, else the usual places), from a running Medorama that holds the
synthetic demo and nothing else: every picture is published. Before it turns Privacy Mode off,
the script checks through the API, signed in as an administrator, that the catalog's only
sources are the demo (\code{DEMO}) and its OMOP database \code{demo\_omop}, and that every
patient and recording came from them; before the first picture, that every patient is one of
the demo's invented patients (each MRN starts with \code{SYN}). It also reads the demo folder's
\file{demo.json}, which \file{make-demo-data.py} writes: it must say that the folder holds
invented data alone, and every recording of the catalog must lie in that folder. A demo made
with \opt{-{}-examples} or \opt{-{}-dicom} holds local ECGs and DICOM studies, so the script
refuses it. If any check fails, it stops without taking a picture. Every request to Chrome has
a time limit.
\begin{params}
\item[OUT] Where the pictures go. Default: \file{docs/images/screens}.
\item[--url URL] The server. Default: \code{http://127.0.0.1:8090/}.
\item[--user NAME] The account. Default: \code{admin}.
\item[--password-file FILE] The file that holds its password. Required.
\item[--demo DIR] The folder the catalog's demo was made in (\file{make-demo-data.py}).
  Required.
\end{params}
\begin{lstlisting}[style=shell]
python3 tools/screenshots.py --url http://127.0.0.1:8090/ --password-file pw --demo /path/to/demo
\end{lstlisting}

\apientry{tools/load-test.py}{Many users at once}
\begin{lstlisting}[style=signature]
python3 tools/load-test.py [--url URL] [--user NAME] [--users N] [--seconds S] [--seed N]
\end{lstlisting}
Simulates users working at once on a Medorama server, each in a session of its own: catalog
pages, patients' pages, database rows and cohort searches, without pause. At the end it prints,
for each kind of request, the count, the median, the 95th percentile, the slowest time and the
errors. It changes nothing (each sign-in is audited, as any). The password is read from the
environment variable \code{MEDORAMA\_PASSWORD}, else asked for.
\begin{params}
\item[--url URL] The server. Default: \code{http://127.0.0.1:8090}.
\item[--user NAME] The account every simulated user signs in as. Default: \code{admin}.
\item[--users N] Users at once. Default: 20.
\item[--seconds S] How long. Default: 60.
\item[--seed N] The seed of the random choices. Default: 1.
\end{params}
\begin{lstlisting}[style=shell]
python3 tools/load-test.py --users 20 --seconds 60
\end{lstlisting}

\apientry{tools/git-hooks/check-files}{Keep patient data out of commits}
\begin{lstlisting}[style=signature]
git config core.hooksPath tools/git-hooks
tools/git-hooks/check-files [--staged | --commits A..B | --tree REV]
\end{lstlisting}
The privacy check. It reads the files as git stores them, never the working tree's copies, and
refuses: anything in a \file{sample\_data/} folder at any depth; recording names of
bedside-monitor systems in a path or in any content, text or binary; files over 5\,MB (this manual's PDF over 20\,MB); notebooks
with saved outputs or that cannot be read; and, outside the synthetic test data and examples,
recordings and signal arrays (WFDB, EDF, \file{.Stp}, NumPy), tables (\file{.csv}, \file{.tsv},
\file{.xlsx}, Parquet), databases and their journals, SQL files and dumps (except the catalog's
schema), archives, and DICOM and GE MUSE XML files, recognized by their content. Without an
option it checks what is staged; \opt{-{}-commits A..B} checks every commit of a range, each
against its parent; \opt{-{}-tree REV} every file of a revision. Its exit status is 0 when
nothing is refused, 1 when something is, 2 on a usage or git error.

The hooks \file{pre-commit} and \file{pre-merge-commit} run it on what a commit or a merge would
store. The first command above enables them once per clone; the installers do it, and warn if
the clone already uses other hooks. CI runs the same check on every commit of each pull request
and push, and on every file (\file{.github/workflows/checks.yml}), so bypassing the hooks
(\opt{-{}-no-verify}) only moves the refusal there.

The \file{Makefile} at the root of the repository wraps these tools and the tests
(\cmd{make help}).
